\documentclass[12pt,a4paper]{report}

% --- Packages ---
\usepackage[utf8]{inputenc}
\usepackage[T1]{fontenc}
\usepackage[english]{babel}
\usepackage{geometry}
\usepackage{fancyhdr}
\usepackage{tcolorbox}
\usepackage{amsmath, amsfonts, amssymb, amsthm}
\usepackage{hyperref}
\usepackage{titlesec}
\usepackage{xcolor}
\usepackage{enumitem}
\usepackage{graphicx}

% --- Page Geometry ---
\geometry{
    a4paper,
    total={170mm,257mm},
    left=20mm,
    top=25mm,
}

% --- Colors ---
\definecolor{mainblue}{RGB}{0, 102, 204}
\definecolor{lightgray}{gray}{0.95}

% --- Hyperref Setup ---
\hypersetup{
    colorlinks=true,
    linkcolor=mainblue,
    filecolor=magenta,      
    urlcolor=cyan,
    pdftitle={Course Notes},
}

% --- Header & Footer ---
\pagestyle{fancy}
\fancyhf{}
\rhead{\leftmark}
\lhead{Course Notes}
\rfoot{Page \thepage}

% --- TColorBox Styles ---
\tcbuselibrary{skins, breakable}

\newtcolorbox{definition}[1][]{
    colback=green!5!white,
    colframe=green!75!black,
    fonttitle=\bfseries,
    title=Definition: #1,
    breakable,
    enhanced,
    attach title to upper,
    after title={:\hspace{0.5em}}
}

\newtcolorbox{theorem}[1][]{
    colback=blue!5!white,
    colframe=blue!75!black,
    fonttitle=\bfseries,
    title=Theorem: #1,
    breakable,
    enhanced,
    attach title to upper,
    after title={:\hspace{0.5em}}
}

\newtcolorbox{example}[1][]{
    colback=orange!5!white,
    colframe=orange!75!black,
    fonttitle=\bfseries,
    title=Example: #1,
    breakable,
    enhanced,
    attach title to upper,
    after title={:\hspace{0.5em}}
}

\newtcolorbox{important}[1][]{
    colback=red!5!white,
    colframe=red!75!black,
    fonttitle=\bfseries,
    title=IMPORTANT,
    breakable,
    enhanced
}

% --- Document Info ---
\title{
    \vspace{2in}
    \Huge \textbf{Multi-Agent Systems} \\
    \Large \textit{Maastricht University} \\
    \vspace{0.5in}
    \large Notes from Lectures \\
    \vspace{2in}
}
\author{Yadder Aceituno}
\date{\today}

% --- Document ---
\begin{document}

\maketitle
\tableofcontents
\newpage

\chapter{Social Choice}

\section{Foundations of Social Choice}
Formally, the issue is combining preferences to derive social outcome.

\begin{definition}
    Agents are called voters and they are the set $Ag = \{1, 2, \dots, n\}$.

    Voters make group decisions that is the set $ \omega = \{w_1, w_2, \dots, w_m\}$ of alternatives.
\end{definition}

If $|\omega| = 2$, we have a pairwise election.

The fundamental problem of social choice is to aggregate the preferences of the voters.
How do we combine the preferences of the voters to make a group decision?

Two main approaches:

\begin{itemize}
    \item Social welfare functions: the output is a social preference order.
    \item Social choice functions: the output is one single alternative.
\end{itemize}

\section{Voting Procedures}

\subsection{Plurality}
Characteristics of this scheme are:

\begin{itemize}
    \item Each voter casts a single vote for their most preferred alternative.
    \item The alternative with the most votes wins.
\end{itemize}

One anomaly of this procedure is for example the outcome of 100 voters is:
\begin{itemize}
    \item 40\% of voters prefer $\omega_1$.
    \item 30\% of voters prefer $\omega_2$.
    \item 30\% of voters prefer $\omega_3$.
\end{itemize}

With plurality rule, $\omega_1$ wins. However, we can see that 60\% of voters do not prefer $\omega_1$.

\end{document}
